\documentclass[10pt,a4paper]{amsart}
\usepackage{amsmath,amssymb,amsthm,graphicx}
% For commutative diagrams we use Xy-pic
\usepackage[all]{xy}

% We use 2cell for 2-commutative diagrams.
\xyoption{2cell}
% Theorem environments.
%
\theoremstyle{plain}
\newtheorem{theorem}[subsection]{Theorem}
\newtheorem{proposition}[subsection]{Proposition}
\newtheorem{lemma}[subsection]{Lemma}

\theoremstyle{definition}
\newtheorem{definition}[subsection]{Definition}
\newtheorem{example}[subsection]{Example}
\newtheorem{exercise}[subsection]{Exercise}
\newtheorem{situation}[subsection]{Situation}

\theoremstyle{remark}
\newtheorem{remark}[subsection]{Remark}
\newtheorem{remarks}[subsection]{Remarks}

\numberwithin{equation}{subsection}

% Macros
%
\def\lim{\mathop{\mathrm{lim}}\nolimits}
\def\colim{\mathop{\mathrm{colim}}\nolimits}
\def\Spec{\mathop{\mathrm{Spec}}}
\def\Hom{\mathop{\mathrm{Hom}}\nolimits}
\def\Ext{\mathop{\mathrm{Ext}}\nolimits}
\def\SheafHom{\mathop{\mathcal{H}\!\mathit{om}}\nolimits}
\def\SheafExt{\mathop{\mathcal{E}\!\mathit{xt}}\nolimits}
\def\Sch{\mathit{Sch}}
\def\Mor{\mathop{\mathrm{Mor}}\nolimits}
\def\Ob{\mathop{\mathrm{Ob}}\nolimits}
\def\Sh{\mathop{\mathit{Sh}}\nolimits}
\def\NL{\mathop{N\!L}\nolimits}
\def\CH{\mathop{\mathrm{CH}}\nolimits}
\def\proetale{{pro\text{-}\acute{e}tale}}
\def\etale{{\acute{e}tale}}
\def\QCoh{\mathit{QCoh}}
\def\Ker{\mathop{\mathrm{Ker}}}
\def\Im{\mathop{\mathrm{Im}}}
\def\Coker{\mathop{\mathrm{Coker}}}
\def\Coim{\mathop{\mathrm{Coim}}}

% Boxtimes
%
\DeclareMathSymbol{\boxtimes}{\mathbin}{AMSa}{"02}

%
% Macros for moduli stacks/spaces
%
\def\QCohstack{\mathcal{QC}\!\mathit{oh}}
\def\Cohstack{\mathcal{C}\!\mathit{oh}}
\def\Spacesstack{\mathcal{S}\!\mathit{paces}}
\def\Quotfunctor{\mathrm{Quot}}
\def\Hilbfunctor{\mathrm{Hilb}}
\def\Curvesstack{\mathcal{C}\!\mathit{urves}}
\def\Polarizedstack{\mathcal{P}\!\mathit{olarized}}
\def\Complexesstack{\mathcal{C}\!\mathit{omplexes}}
% \Pic is the operator that assigns to X its picard group, usage \Pic(X)
% \Picardstack_{X/B} denotes the Picard stack of X over B
% \Picardfunctor_{X/B} denotes the Picard functor of X over B
\def\Pic{\mathop{\mathrm{Pic}}\nolimits}
\def\Picardstack{\mathcal{P}\!\mathit{ic}}
\def\Picardfunctor{\mathrm{Pic}}
\def\Deformationcategory{\mathcal{D}\!\mathit{ef}}
\newenvironment{reference}{}{}
\newenvironment{slogan}{}{}
\newenvironment{history}{}{}
\usepackage{cleveref}
\crefname{originallemma}{original lemma}{original lemmas}
\Crefname{originallemma}{Original lemma}{Original lemmas}
\makeatletter
\newlabel{lemma-bound-neighbours}{{0C9U}{source}{}{sourcepointer-1}{}}
\newlabel{lemma-bound-neighbours@cref}{{[originallemma][0][]0C9U}{[1][0][]source}{}{}{}}
\makeatother
\crefname{originaldefinition}{original definition}{original definitions}
\Crefname{originaldefinition}{Original definition}{Original definitions}
\makeatletter
\newlabel{definition-type-minimal}{{0C7A}{source}{}{sourcepointer-2}{}}
\newlabel{definition-type-minimal@cref}{{[originaldefinition][0][]0C7A}{[1][0][]source}{}{}{}}
\makeatother
\crefname{originaldefinition}{original definition}{original definitions}
\Crefname{originaldefinition}{Original definition}{Original definitions}
\makeatletter
\newlabel{definition-type-minus-one}{{0C76}{source}{}{sourcepointer-3}{}}
\newlabel{definition-type-minus-one@cref}{{[originaldefinition][0][]0C76}{[1][0][]source}{}{}{}}
\makeatother
\crefname{originaldefinition}{original definition}{original definitions}
\Crefname{originaldefinition}{Original definition}{Original definitions}
\makeatletter
\newlabel{definition-type-minus-two}{{0C7E}{source}{}{sourcepointer-4}{}}
\newlabel{definition-type-minus-two@cref}{{[originaldefinition][0][]0C7E}{[1][0][]source}{}{}{}}
\makeatother
\crefname{originallemma}{original lemma}{original lemmas}
\Crefname{originallemma}{Original lemma}{Original lemmas}
\makeatletter
\newlabel{lemma-diagonal-negative}{{0C74}{source}{}{sourcepointer-5}{}}
\newlabel{lemma-diagonal-negative@cref}{{[originallemma][0][]0C74}{[1][0][]source}{}{}{}}
\makeatother
\crefname{originallemma}{original lemma}{original lemmas}
\Crefname{originallemma}{Original lemma}{Original lemmas}
\makeatletter
\newlabel{lemma-minus-one}{{0C75}{source}{}{sourcepointer-6}{}}
\newlabel{lemma-minus-one@cref}{{[originallemma][0][]0C75}{[1][0][]source}{}{}{}}
\makeatother
\crefname{originallemma}{original lemma}{original lemmas}
\Crefname{originallemma}{Original lemma}{Original lemmas}
\makeatletter
\newlabel{lemma-minus-two}{{0C7D}{source}{}{sourcepointer-7}{}}
\newlabel{lemma-minus-two@cref}{{[originallemma][0][]0C7D}{[1][0][]source}{}{}{}}
\makeatother
\title{A bound on the heart of a minimal numerical type of genus at least two}
\author{The Stacks Project Authors}
\date{stacks-project:models.tex:local-snapshot-20261012; GNU FDL 1.2 or later}
\begin{document}
\maketitle
\noindent\textit{Editorial note.} The counted unit is the Stacks Project Lemma lemma-bound-heart at source lines 4114-4126 together with its proof at source lines 4128-4169, a straight adjacent pairing in the chapter file. The chapter is the upstream file models.tex; the pool directory carries the longer name stacks-models and is not used as the chapter name anywhere in this card.\par\medskip
\noindent\textit{Editorial note.} This card has no context span: the statement and proof are self-contained apart from the seven pointers, and the notions they use are recorded in the standard prerequisites.\par\medskip
\noindent\textit{Editorial note.} The retained proof names seven results of the same chapter: an elementary bound, a definition of minimal types, two boundary definitions and three lemmas. All are declared as non-load-bearing navigation pointers with their stable Stacks tags looked up in the shared tag index. Two of them belong to the block whose proofs this lane does not count, because those proofs are one-line citations; pointing at their statements is nevertheless valid and none of their intervals is reproduced here.\par\medskip
\noindent\textit{Editorial note.} The card on the weighted intersection bound at lines 4171-4283, registered earlier by this lane, points at this lemma; the pointers and this counted unit are therefore consistent.\par\medskip
\noindent\textit{Editorial note.} The retained proof span starts at line 4128, the line of the opening marker read from the source, and contains exactly one proof environment.\par\medskip
\noindent\textit{Editorial note.} No citation occurs in the retained spans, so no bibliography entry is supplied.\par\medskip
\noindent\textit{Allowed standard prerequisites (editorial).}\begin{itemize}
\item Minimal numerical types, minus-one and minus-two indices and the elementary bounds. A numerical type is a tuple of an integer and integer families of multiplicities, intersection numbers, weights and genera subject to the five defining conditions of the chapter; it is minimal when it has no minus-one index, that is, no index whose genus is zero and whose self-intersection equals the negative of its weight, and the boundary cases where the diagonal entry is large are collected in the definitions of minus-one and minus-two indices. The chapter's elementary bound on neighbouring entries and its lemmas on the sign of the diagonal entries and on the genus of a neighbouring index supply the estimates the proof uses. The author's own definitions and lemmas of the same chapter are the allowed inputs.. The counted lemma is the quantitative input of the classification of minimal types: together with the bound on the weighted intersection numbers it makes the heart of the type bounded, which is what makes the enumeration of minimal types of fixed genus finite.
\end{itemize}\medskip
\ifdefined\hypertarget\hypertarget{sourcepointer-1}{}\fi
\noindent\textit{Source reference.} Original lemma 0C9U: focused-seeds/AGP/sources/stacks-snapshot-20261011/tags.txt, source lines 15414--15414. Navigation pointer invoked by the retained proof at its first line to reduce the assertion to a single entry. This lane does not count that lemma, because its own proof defers the argument to the reader; pointing at its statement from another card is nevertheless valid and its interval is not reproduced here. The number is the stable Stacks tag of the target, looked up in the shared tag index.\par\medskip
\ifdefined\hypertarget\hypertarget{sourcepointer-2}{}\fi
\noindent\textit{Source reference.} Original definition 0C7A: focused-seeds/AGP/sources/stacks-snapshot-20261011/tags.txt, source lines 15325--15325. Navigation pointer to the chapter's definition of a minimal numerical type, used by the retained proof to exclude the minus-one configurations; the cited definition is an allowed input and its text is not reproduced here. The number is the stable Stacks tag of the target, looked up in the shared tag index.\par\medskip
\ifdefined\hypertarget\hypertarget{sourcepointer-3}{}\fi
\noindent\textit{Source reference.} Original definition 0C76: focused-seeds/AGP/sources/stacks-snapshot-20261011/tags.txt, source lines 15321--15321. Navigation pointer to the chapter's definition of the boundary configuration with a large diagonal entry; the cited definition is an allowed input and its text is not reproduced here. The number is the stable Stacks tag of the target, looked up in the shared tag index.\par\medskip
\ifdefined\hypertarget\hypertarget{sourcepointer-4}{}\fi
\noindent\textit{Source reference.} Original definition 0C7E: focused-seeds/AGP/sources/stacks-snapshot-20261011/tags.txt, source lines 15329--15329. Navigation pointer to the chapter's second boundary definition, used by the retained proof in its case analysis; the cited definition is an allowed input and its text is not reproduced here. The number is the stable Stacks tag of the target, looked up in the shared tag index.\par\medskip
\ifdefined\hypertarget\hypertarget{sourcepointer-5}{}\fi
\noindent\textit{Source reference.} Original lemma 0C74: focused-seeds/AGP/sources/stacks-snapshot-20261011/tags.txt, source lines 15319--15319. Navigation pointer invoked by the retained proof for the sign of the diagonal entries of a type with more than one vertex. This lane does not count that lemma because its own proof is a bare citation; pointing at its statement is nevertheless valid and its interval is not reproduced here. The number is the stable Stacks tag of the target, looked up in the shared tag index.\par\medskip
\ifdefined\hypertarget\hypertarget{sourcepointer-6}{}\fi
\noindent\textit{Source reference.} Original lemma 0C75: focused-seeds/AGP/sources/stacks-snapshot-20261011/tags.txt, source lines 15320--15320. Navigation pointer invoked by the retained proof in the analysis of the first boundary configuration. This lane does not count that lemma because its own proof is a bare citation; pointing at its statement is nevertheless valid and its interval is not reproduced here. The number is the stable Stacks tag of the target, looked up in the shared tag index.\par\medskip
\ifdefined\hypertarget\hypertarget{sourcepointer-7}{}\fi
\noindent\textit{Source reference.} Original lemma 0C7D: focused-seeds/AGP/sources/stacks-snapshot-20261011/tags.txt, source lines 15328--15328. Navigation pointer invoked by the retained proof in the analysis of the second boundary configuration; the cited statement is an allowed input of the same chapter and its proof is not reproduced here. The number is the stable Stacks tag of the target, looked up in the shared tag index.\par\medskip
\begin{lemma}
\label{lemma-bound-heart}
Fix $g \geq 2$. For every minimal numerical type $n, m_i, a_{ij}, w_i, g_i$
of genus $g$ with $n > 1$ we have
\begin{enumerate}
\item the set $J \subset \{1, \ldots, n\}$ of non-$(-2)$-indices
has at most $2g - 2$ elements,
\item for $j \in J$ we have $g_j < g$,
\item for $j \in J$ we have $m_j|a_{jj}| \leq 6g - 6$, and
\item for $j \in J$ and $i \in \{1, \ldots, n\}$
we have $m_ia_{ij} \leq 6g - 6$.
\end{enumerate}
\end{lemma}

\begin{proof}
Recall that $g = 1 + \sum m_j(w_j(g_j - 1) - \frac{1}{2} a_{jj})$.
For $j \in J$ the contribution $m_j(w_j(g_j - 1) - \frac{1}{2} a_{jj})$
to the genus $g$ is $> 0$ and hence $\geq 1/2$. This uses
Lemma \ref{lemma-minus-one},
Definition \ref{definition-type-minus-one},
Definition \ref{definition-type-minimal},
Lemma \ref{lemma-minus-two}, and
Definition \ref{definition-type-minus-two}; we will use these
results without further mention in the following.
Thus $J$ has at most $2(g - 1)$ elements.
This proves (1).

\medskip\noindent
Recall that $-a_{ii} > 0$ for all $i$ by Lemma \ref{lemma-diagonal-negative}.
Hence for $j \in J$ the contribution $m_j(w_j(g_j - 1) - \frac{1}{2} a_{jj})$
to the genus $g$ is $> m_jw_j(g_j - 1)$. Thus
$$
g - 1 > m_jw_j(g_j - 1) \Rightarrow g_j < (g - 1)/m_jw_j + 1
$$
This indeed implies $g_j < g$ which proves (2).

\medskip\noindent
For $j \in J$ if $g_j > 0$, then the contribution
$m_j(w_j(g_j - 1) - \frac{1}{2} a_{jj})$ to the genus $g$
is $\geq -\frac{1}{2}m_ja_{jj}$ and we immediately conclude
that $m_j|a_{jj}| \leq 2(g - 1)$. Otherwise $a_{jj} = -kw_j$
for some integer $k \geq 3$ (because $j \in J$) and we get
$$
m_jw_j(-1 + \frac{k}{2}) \leq g - 1
\Rightarrow
m_jw_j \leq \frac{2(g - 1)}{k - 2}
$$
Plugging this back into $a_{jj} = -km_jw_j$ we obtain
$$
m_j|a_{jj}| \leq 2(g - 1) \frac{k}{k - 2} \leq 6(g - 1)
$$
This proves (3).

\medskip\noindent
Part (4) follows from Lemma \ref{lemma-bound-neighbours} and (3).
\end{proof}

\end{document}
